% OWNER: Lynn (L6). One paragraph, ~150-200 words.
\begin{abstract}
Large language models are increasingly considered as planning components, yet on the open-ended
questions of ACPBench-Hard even large models remain far from reliable. We reproduce the benchmark's
Next Action task, on which a model must generate an action that moves a state one optimal step
closer to the goal; with the Progression task, we match all four reported results within their 95\% intervals. We then ask whether LLM-guided best-first search, in which the model
proposes actions, predicts their effects and scores the predicted states without any planner or
verifier, improves on direct generation under a matched call budget. It does not: with
Llama~3.1~8B, search is lower than direct generation in every seed, though not significantly ($0.159$ vs.\ $0.223$), and with
70B not significantly different ($0.362$ vs.\ $0.377$), while on 8B, best-of-$N$ with the same budget scores higher than search. The model often proposes a correct action (in 31\% of questions for 8B and 59\% for
70B) but often fails to recognise it: 8B ranks its candidates no better than chance, even 70B picks the correct one in only 61\% of those questions, and no evaluator we tried on 8B beats trusting the proposer's own order. A post-hoc follow-up that never
judges its own candidates, self-consistency voting over structured reasoning, reaches $0.331$ on 8B and $0.492$ on 70B, above best-of-$N$ with the same 14 calls ($0.231$, $0.364$; significant after Holm correction), though its gain over one greedy structured call ($0.262$, $0.469$) is not significant. For open-ended planning without
a verifier, evaluation, more than generation, is the bottleneck.
\end{abstract}
